% Chapter 9: Peak
% Scene function: The haunting at technical peak. Tert's best work.
% Gordon laughs at 4 AM. The margin pad stays closed.
% Pre-echo of the turn.

\chapter{Wednesday Night}

Week three. The haunting is at technical peak.

I have been in this house for seventeen days and I know it the way a musician knows an instrument. I know which joist creaks at what interval. I know how the pipes resonate against the foundation walls. I know that the attic ladder, if released by half an inch from its catch, produces a sound at 2 a.m. that is indistinguishable from a footstep on the landing. I know where to stand in the hallway to be visible in the bathroom mirror for exactly one-quarter second when the subject passes. I know the house's thermal rhythms, its air currents, the way sound travels from the basement to the bedroom through the gap under the sickroom door.

I have not used the gap under the sickroom door. Some things you hold in reserve.

Gordon is not sleeping well. This is evident in everything: the pace of his footsteps in the morning, the length of his showers, the number of times he reheats the same cup of coffee before giving up on it. He has started leaving lights on in rooms he is not in, which is a Phase 2 indicator so classical it is literally in the manual. A subject who leaves lights on is a subject whose relationship with darkness has shifted. They may not know why. They do not need to.

On Wednesday night I produce my best work of the assignment.

It is a sequence. Not a single event. A three-part escalation over ninety minutes, using the bookshelf wall, the bathroom tap, and a thermal pocket I have been cultivating in the hallway outside the bedroom since week one. The sequence is designed so that each element arrives just as the previous one has faded below threshold: the creak, then silence; the tap, then silence; then the cold spot in the hallway, timed to coincide with Gordon's first REM cycle, so that when his body registers the temperature drop it incorporates the sensation into whatever he is dreaming and the dream turns.

This is atmospheric work at the level I have always known I was capable of and have rarely had a canvas for. The timing is precise. The layering is clean. The house performs exactly as I have taught it to perform. If the departmental newsletter existed, this would make the cover.

Gordon wakes at 3:47. He does not get out of bed. He lies still for several minutes. I can hear his breathing from the hallway: elevated, then deliberate, then controlled. He is calming himself down. A subject who has to work to calm down is a subject whose baseline has been successfully lowered.

He gets up at four and goes to the kitchen. He does not turn on the hall light. I note this. A subject who was leaving lights on but now walks through darkness has either habituated or entered a state of fatigue where the effort of reaching for the switch exceeds the comfort of the light. Both are useful.

He makes coffee. It is four in the morning and he is making coffee because I woke him and he has decided that going back to sleep is not something he can do right now. He stands at the counter in the dark kitchen with the coffee maker running and the only light coming from the microwave display, and he looks out the window at the back yard where the garden is.

And he laughs.

Not loudly. Not at anything I can identify. A short sound, almost a breath, that starts as something else and becomes a laugh halfway through. He shakes his head once. He is laughing at himself, or at the situation, or at something I am not equipped to observe.

The laugh does not fit the ambient environment I have spent three weeks building. It is evidence of something inside Gordon that my work has not touched.

I stand in the hallway and I listen to him pour his coffee and I do not make a note in my margin pad.
